% Conclusion.
%
% Says what was found and what it licenses, and nothing that is not already
% established in Sections 5 and 6. No number appears here that does not appear
% in a generated table.

\section{Conclusion}
\label{sec:conclusion}

We set out to build the missing instrument --- a benchmark for catastrophic
forgetting in the transition component of a world model, rather than in a
policy or in an agent as a whole --- and the instrument's most useful output was
a correction to two things the field does by default.

The first is a design habit. Continual learning experiments are routinely built
around an ordinal notion of how far apart the tasks in a sequence are, and that
ordering is then treated as an experimental factor. Ours does not order
forgetting: RD peaks at the medium level in all three families, the labelled
level carries no rank information across the nine cells, and the cleanest case
needs no comparison between families at all --- one family's maximum level is
its medium perturbation plus two more, on the same task pair, and it forgets
less. That survives a doubling of the training budget, a doubling of the seed
count, and a four-task sequence, and in the one case we could interrogate
directly the obvious explanation turns out to be backwards: the more heavily
perturbed task is the \emph{easier} one to fit. What we can say is that
forgetting tracked what the new task demanded of the model rather than how far
it sat from the old one. What we cannot yet say is which of the two measurable
quantities standing in for that demand should be reported, since they swapped
rank when we added seeds.

The second is a measurement habit. Isolating a component means holding a
representation fixed, and a metric defined on a fixed representation cannot see
that representation move. Ours cannot: fine-tuning's held-out reconstruction of
the first task degrades by a factor of $811$ while its prediction fidelity
records an improvement. EWC turns that from an artefact into a structural
prediction --- its Fisher information is identically zero on every encoder
parameter, so it protects the transition component almost exactly and nothing
else --- and both halves of the prediction were checkable before the experiment
ran. A component-level benchmark owes its readers both scales, and ours reports
them side by side for that reason.

Neither result is about which method wins, and neither is comfortable for the
method we proposed: UG-MTM does not forget because it freezes its encoder, and
pays for it in transfer. We report it as one of five characterised methods.

The limits are the ones stated in Section~\ref{sec:results:threats} and we do
not think they should be read past. This is one architecture family at a small
budget on data from a random policy, and nothing here establishes what happens
at the scale world models are actually trained at. The natural next step is not
a better aggregate or another mitigation method but the same measurements on a
model somebody would deploy --- and the second is a distance between
environments that survives contact with its own instrument, which ours does
not.

Code, configurations, the plotting and table-generating scripts, and all 375
result files are released.
